\documentclass[12pt,a4paper]{article}

% ─── Packages ────────────────────────────────────────────────────────────────
\usepackage[utf8]{inputenc}
\usepackage[T1]{fontenc}
\usepackage{geometry}
\geometry{margin=2.5cm}

\usepackage{amsmath, amssymb}
\usepackage{algorithm}
\usepackage{algpseudocode}
\usepackage{booktabs}
\usepackage{array}
\usepackage{xcolor}
\usepackage{enumitem}
\usepackage{tcolorbox}
\tcbuselibrary{skins, breakable}
\usepackage{hyperref}
\usepackage{tikz}
\usetikzlibrary{arrows.meta, positioning, shapes.geometric}

% ─── Colors ──────────────────────────────────────────────────────────────────
\definecolor{urgence}{RGB}{180,30,30}
\definecolor{theorie}{RGB}{30,80,160}
\definecolor{pratique}{RGB}{20,130,80}
\definecolor{attention}{RGB}{200,120,0}
\definecolor{bleuléger}{RGB}{230,240,255}
\definecolor{rougeclair}{RGB}{255,235,235}
\definecolor{orangeclair}{RGB}{255,245,225}
\definecolor{grisléger}{RGB}{245,245,248}
\definecolor{noirbloc}{RGB}{25,25,40}

% ─── Custom boxes ─────────────────────────────────────────────────────────────
\newtcolorbox{defbox}[1]{
  colback=bleuléger, colframe=theorie, fonttitle=\bfseries,
  title={#1}, breakable, sharp corners=south
}
\newtcolorbox{alertbox}[1]{
  colback=rougeclair, colframe=urgence, fonttitle=\bfseries,
  title={#1}, breakable, sharp corners=south
}
\newtcolorbox{methbox}[1]{
  colback=grisléger, colframe=gray!50, fonttitle=\bfseries,
  title={#1}, breakable, sharp corners=south
}
\newtcolorbox{exemplebox}[1]{
  colback=orangeclair, colframe=attention, fonttitle=\bfseries,
  title={#1}, breakable, sharp corners=south
}

% ─── Hyperref ─────────────────────────────────────────────────────────────────
\hypersetup{
  colorlinks=true,
  linkcolor=theorie,
  urlcolor=urgence,
  pdftitle={Twist 4 — Routes fermées brutalement},
  pdfauthor={Songo — Simulation de routage},
}

% ─────────────────────────────────────────────────────────────────────────────
\begin{document}

\begin{center}
  {\LARGE\bfseries Twist 4 : Routes ferm\'ees brutalement}\\[6pt]
  {\large Recalcul en cours d'intervention et justification de d\'ecision}\\[3pt]
  {\small\color{gray} Songo MVP --- Routage dynamique A*}
\end{center}

\vspace{0.4cm}

\noindent
Dans une intervention r\'eelle, l'environnement ne reste pas fig\'e.
Un axe parfaitement praticable au d\'epart peut \^etre \textbf{ferm\'e sans pr\'evenir}
en cours de route : accident, foule, barrage de s\'ecurit\'e, effondrement, inondation soudaine.
Le syst\`eme doit alors r\'eagir \textbf{imm\'ediatement}, recalculer un nouvel
itin\'eraire depuis la position courante du v\'ehicule, et \textbf{justifier
le changement} --- tant pour le conducteur que pour le jury.

\bigskip

% ── 1. D\'efinition du probl\`eme ────────────────────────────────────────────
\begin{defbox}{D\'efinition --- Fermeture brutale d'un arc}
Soit $G = (V, E, w)$ le graphe routier dynamique, o\`u $w(e, t)$ est le poids
de l'arc $e$ au temps simul\'e $t$.
Une \textbf{fermeture brutale} est la transition instantan\'ee~:
\[
  w(e,\, t^-) \;<\; +\infty
  \quad\longrightarrow\quad
  w(e,\, t^+) = +\infty
\]
\`a un instant $t$ imprévisible, alors que le v\'ehicule est d\'ej\`a en route.
L'arc $e$ est ainsi retir\'e de l'espace de recherche pour tout recalcul ult\'erieur.
\end{defbox}

\bigskip

% ── 2. Pourquoi c'est critique ───────────────────────────────────────────────
\begin{alertbox}{Pourquoi ce twist est critique pour un v\'ehicule d'urgence}
\begin{itemize}[leftmargin=1.5em, itemsep=4pt]
  \item \textbf{L'itin\'eraire initial devient invalide} : le v\'ehicule suit
        un chemin qui n'existe plus ; continuer sans recalcul peut mener
        \`a un cul-de-sac ou \`a un danger direct.
  \item \textbf{Le temps perdu est irr\'ecup\'erable} : chaque seconde de
        h\'esitation devant une route bloqu\'ee aggrave la situation
        \`a destination (patient, incendie, incident).
  \item \textbf{La position courante change le probl\`eme} : le v\'ehicule
        n'est plus \`a son point de d\'epart initial ; le recalcul doit partir
        de l\`a o\`u il se trouve \textit{maintenant}, pas du d\'epart original.
  \item \textbf{La justification est obligatoire} : le conducteur doit comprendre
        pourquoi l'itin\'eraire change, et le jury doit voir que la d\'ecision
        est logique et tra\c{c}able.
\end{itemize}
\end{alertbox}

\bigskip

% ── 3. M\'ecanisme de d\'etection ────────────────────────────────────────────
\section*{1. D\'etection de la fermeture}

Trois sources peuvent d\'eclencher une fermeture brutale dans la simulation~:

\begin{enumerate}[leftmargin=1.5em, itemsep=4pt]
  \item \textbf{\'Ev\'enement manuel} : le d\'emonstrateur clique sur le bouton
        ``Bloquer axe'' dans l'interface ; un marqueur rouge appara\^it sur la carte
        et l'arc correspondant est mis \`a $+\infty$ imm\'ediatement.
  \item \textbf{\'Ev\'enement automatiqu\'e} : le moteur injecte automatiquement
        un blocage \`a $t + 20$~minutes simul\'ees, sur un arc situ\'e
        en avant du v\'ehicule, pour illustrer le sc\'enario sans intervention
        humaine.
  \item \textbf{\'Ev\'enement contextuel} : une pluie soudaine ($t + 10$~min)
        peut rendre certains axes impraticables ($w \to +\infty$), ce qui
        constitue une fermeture fonctionnelle m\^eme sans blocage physique.
\end{enumerate}

\noindent
Dans tous les cas, d\`es que la mise \`a jour du graphe est d\'etect\'ee,
le moteur v\'erifie si \textbf{l'arc bloqu\'e appartient au chemin courant}.
Si oui, le recalcul est d\'eclench\'e sans d\'elai.

\bigskip

% ── 4. Proc\'edure de recalcul ───────────────────────────────────────────────
\section*{2. Proc\'edure de recalcul depuis la position courante}

\begin{methbox}{Algorithme de r\'eponse \`a une fermeture}
\begin{enumerate}[leftmargin=1.5em, itemsep=3pt]
  \item \textbf{Gel de la position} : la position exacte $p_{\text{now}}$ du
        v\'ehicule est captur\'ee par interpolation sur le chemin parcouru.
  \item \textbf{Mise \`a jour du graphe} : l'arc ferm\'e re\c{c}oit le poids
        $w = +\infty$ ; les noeuds de d\'etour potentiels conservent leurs
        poids courants (pluie, congestion, heure).
  \item \textbf{Relance de A*} : l'algorithme repart de $p_{\text{now}}$
        vers la destination finale, avec le graphe mis \`a jour.
        Le chemin d\'ej\`a parcouru est \textbf{conserv\'e} et affich\'e
        en pointill\'es sur la carte ; seul le tron\c{c}on restant est recalcul\'e.
  \item \textbf{Validation du r\'esultat} : si A* trouve un chemin alternatif,
        il remplace imm\'ediatement le tron\c{c}on restant.
        Sinon, un avertissement est \'emis (``aucune alternative'').
  \item \textbf{Journalisation} : le moteur enregistre dans le journal la raison
        du recalcul, la position au moment du blocage, la diff\'erence de distance,
        et le nouveau temps estim\'e.
\end{enumerate}
\end{methbox}

\bigskip

\noindent
Le pseudo-code ci-dessous formalise la r\'eponse \`a un \'ev\'enement de fermeture~:

\begin{algorithm}[H]
\caption{R\'eponse \`a une fermeture d'arc en cours d'intervention}
\begin{algorithmic}[1]
\Require position courante $p$, destination $d$, graphe $G$, arc ferm\'e $e^*$
\State $G \leftarrow G$ avec $w(e^*) \leftarrow +\infty$
\If{$e^* \in \text{chemin\_restant}(p, d)$}
    \State Suspendre l'avancement du v\'ehicule
    \State $\pi_{\text{new}} \leftarrow \text{A*}(G,\; p,\; d)$
    \If{$\pi_{\text{new}} \neq \emptyset$}
        \State Remplacer le tron\c{c}on restant par $\pi_{\text{new}}$
        \State Journaliser(raison, $\Delta$km, $\Delta$temps, arc bloqu\'e)
        \State Reprendre l'avancement sur $\pi_{\text{new}}$
    \Else
        \State \'Emettre avertissement ``aucune alternative trouv\'ee''
    \EndIf
\Else
    \State \textit{L'arc ferm\'e n'est pas sur le chemin --- aucune action}
\EndIf
\end{algorithmic}
\end{algorithm}

\bigskip

% ── 5. Justification de la d\'ecision ────────────────────────────────────────
\section*{3. Justification de la d\'ecision}

Le syst\`eme ne se contente pas de recalculer en silence.
Chaque recalcul produit une \textbf{entr\'ee dans le journal de simulation}
visible dans l'interface, structur\'ee comme suit~:

\begin{center}
\begin{tabular}{lp{9cm}}
\toprule
\textbf{Champ} & \textbf{Contenu} \\
\midrule
Cause          & Nature de l'\'ev\'enement (blocage manuel, pluie, fermeture auto) \\
Instant        & Heure simul\'ee au moment du blocage \\
Position       & Fraction du parcours ($\%$) et kilom\'etrage parcouru \\
Ancien chemin  & Distance restante avant blocage \\
Nouveau chemin & Distance et dur\'ee estim\'ee de l'alternative \\
Gain / Perte   & Diff\'erence en km et en secondes par rapport \`a l'ancien chemin \\
Qualit\'e alt. & Signal si la d\'eviation est significativement diff\'erente \\
\bottomrule
\end{tabular}
\end{center}

\medskip
\noindent
Cette tra\c{c}abilit\'e remplit deux r\^oles~: elle informe le conducteur
en temps r\'eel, et elle d\'emontre au jury que le moteur \textbf{ne choisit
pas au hasard} --- chaque d\'eviation est motiv\'ee par des donn\'ees objectives.

\bigskip

% ── 6. Exemple concret ───────────────────────────────────────────────────────
\begin{exemplebox}{Exemple simul\'e --- Blocage \`a $t + 20$ minutes}
\textbf{Contexte} : ambulance en route depuis le quartier Bastos vers l'H\^opital
Central. \`A $t + 20$~minutes simul\'ees, l'avenue Kennedy est ferm\'ee
(accident simulé).

\medskip
\begin{tabular}{ll}
\toprule
\textbf{\'Etape} & \textbf{\'Etat du syst\`eme} \\
\midrule
$t = 19$~min  & Ambulance \`a 42\,\% du parcours, sur l'avenue Kennedy \\
$t = 20$~min  & Fermeture d\'etect\'ee --- $w(\text{Kennedy}) \leftarrow +\infty$ \\
$t = 20$~min  & A* relanc\'e depuis la position courante \\
$t = 20$~min  & Alternative : Boulevard de l'URSS $+$ rue de la Joie \\
Journal       & \textit{``Blocage auto \`a t+20 min --- contournement, +0.8 km, +42 s''} \\
Affichage     & Ancien tron\c{c}on en pointill\'es, nouveau en rouge plein \\
\bottomrule
\end{tabular}

\medskip
\noindent
Le conducteur voit imm\'ediatement le changement sur la carte et lit la raison
dans le journal.
Le jury constate que le moteur a identifi\'e la fermeture, quantifi\'e
le surco\^ut, et propos\'e une alternative coh\'erente.
\end{exemplebox}

\bigskip

% ── 7. Comportement selon le type de v\'ehicule ──────────────────────────────
\section*{4. Comportement selon le type de v\'ehicule}

La r\'eponse \`a une fermeture est identique pour les trois v\'ehicules,
mais le \textbf{choix de l'alternative} diff\`ere selon leurs contraintes~:

\medskip
\begin{tabular}{lp{9.5cm}}
\toprule
\textbf{V\'ehicule} & \textbf{Comportement lors du recalcul apr\`es fermeture} \\
\midrule
Ambulance
  & Cherche l'alternative la plus rapide, m\^eme si elle passe par
    des axes \'etroits ou d\'egrad\'es.
    Priorit\'e absolue au temps. \\[4pt]
Pompiers
  & L'alternative doit satisfaire les contraintes de gabarit.
    Une ruelle trop \'etroite ou une piste en terre sera ignor\'ee,
    m\^eme si elle est plus courte. \\[4pt]
Police
  & Compromis : l'alternative favorise les grands axes mais accepte
    une voie secondaire en bon \'etat si elle reste praticable. \\
\bottomrule
\end{tabular}

\medskip
\noindent
Ainsi, face au \textbf{m\^eme blocage}, les trois v\'ehicules peuvent
emprunter des itin\'eraires diff\'erents --- ce qui illustre directement
l'int\'er\^et de la diff\'erenciation des profils pr\'esent\'ee dans
le document pr\'ec\'edent.

\bigskip

% ── Conclusion ───────────────────────────────────────────────────────────────
\begin{tcolorbox}[colback=grisléger, colframe=theorie, fonttitle=\bfseries,
  title={En r\'esum\'e}, sharp corners=south]
Le Twist~4 met en \'evidence la capacit\'e du moteur \`a g\'erer
l'\textbf{impr\'evisible en temps r\'eel}.
Trois propri\'et\'es cl\'es le caract\'erisent~:
\begin{itemize}[leftmargin=1.5em, itemsep=3pt]
  \item \textbf{R\'eactivit\'e} : le recalcul est instantan\'e d\`es la d\'etection
        de la fermeture, sans intervention humaine.
  \item \textbf{Coh\'erence} : le chemin d\'ej\`a parcouru est pr\'eserv\'e ;
        seul le tron\c{c}on restant est recalcul\'e depuis la position exacte.
  \item \textbf{Tra\c{c}abilit\'e} : chaque d\'ecision est journalis\'ee avec
        sa cause, ses co\^uts et la comparaison avec l'ancien itin\'eraire.
\end{itemize}
\end{tcolorbox}

\end{document}